\uuid{GiUJ}
\exo7id{205}
\titre{exo7 205}
\auteur{ridde}
\organisation{exo7}
\datecreate{1999-11-01}
\isIndication{false}
\isCorrection{false}
\chapitre{Injection, surjection, bijection}
\sousChapitre{Bijection}
\cpge{mpsi-mp2i.raisonnement-vocabulaire-ensembliste.c, pcsi.raisonnement-vocabulaire-ensembliste.d, ptsi.raisonnement-vocabulaire-ensembliste.d, mpsi-mp2i.raisonnement-vocabulaire-ensembliste.b, pcsi.raisonnement-vocabulaire-ensembliste.b, ptsi.raisonnement-vocabulaire-ensembliste.b}
\module{Algèbre}
\niveau{L1}
\difficulte{}

\contenu{
\question{
Soit $X$ un ensemble. Si $A \subset X$ on note $\chi _A$ la fonction caract\'eristique
associ\'ee. Montrer que ${\Phi} : {\mathcal{P} (X)} \to {\mathcal{F} (X, \left\{ 0, 1 \right\})}$, ${A} \mapsto {\chi _A}$ est bijective.
}
}
